\chapter{Sperimentazione}
\label{ch:sperimentazione}

Il presente capitolo è dedicato alla validazione empirica dei principi teorici esaminati nei capitoli precedenti. Attraverso una suite di test controllati, verranno analizzati i comportamenti dinamici del database graph-native Neo4j e del sistema relazionale PostgreSQL, ponendo particolare enfasi sui tempi di risposta, sulla gestione della concorrenza transazionale e sulla resilienza in contesti distribuiti.

\section{Setup dell'Ambiente e Metodologia di Test}
\label{sec:setup_metodologia}

\subsection{Il Dataset di Riferimento}
Per garantire l'oggettività e la riproducibilità dei benchmark, la sperimentazione adotta l'\textbf{LDBC Social Network Benchmark (SNB)} nello scenario \textit{Interactive}. Si tratta di un benchmark standard industriale progettato dal Linked Data Benchmark Council per simulare i carichi di lavoro tipici di un social network.

Il \textit{Datagen} di LDBC produce un insieme di file CSV strutturati e normalizzati che riflettono un dominio fortemente interconnesso (entità \texttt{Person}, \texttt{Post}, \texttt{Comment}, \texttt{Forum} e \texttt{Tag}). Questo approccio permette di popolare entrambi i sistemi partendo dalla medesima base informativa tabellare, garantendo la perfetta simmetria dei dati estratti.

\subsection{Architettura Hardware e Macchine Virtuali}
L'infrastruttura di test è composta da un cluster di macchine virtuali basate su architettura Linux Ubuntu Server, allocate all'interno dell'ambiente di laboratorio della Scuola di Scienze. Ciascuna istanza dispone di specifiche dedicate:
\begin{itemize}
  \item \textbf{CPU:} 4 Core dedicati vCPU.
  \item \textbf{RAM:} 8 GB di memoria volatile.
  \item \textbf{Storage:} Memoria a stato solido SSD ad alte prestazioni.
  \item \textbf{Rete:} Connessione inter-nodo privata a bassa latenza.
\end{itemize}
PostgreSQL è configurato in un'istanza singola ottimizzata (allocazione della \textit{shared\_buffers} pari al 25\% della RAM totale). Neo4j è distribuito in una configurazione flessibile tramite container Docker, orchestrati sia in modalità a nodo singolo (per gli scenari 1, 2 e 4) sia in un cluster Raft distribuito a 5 nodi per l'analisi dei sistemi distribuiti.

\subsection{Metriche di Valutazione}
I carichi di lavoro sono orchestrati tramite uno script di automazione sviluppato in Python, che interagisce con i database mediante le librerie ufficiali \texttt{psycopg2} e \texttt{neo4j}. Le metriche quantitative monitorate sono:
\begin{itemize}
  \item \textbf{Latenza di Esecuzione:} Espressa in millisecondi (ms), calcolata registrando il tempo CPU ad alta precisione tramite la funzione \texttt{time.perf\_counter()}. Per ogni query viene calcolata la media e il 90\textsuperscript{o} percentile su un campione di ripetizioni iterative.
  \item \textbf{Throughput:} Misurato in query per secondo (QPS) sotto stress test massivi.
  \item \textbf{Saturazione delle Risorse:} Monitoraggio della memoria RAM allocata e del carico percentuale della CPU durante le computazioni critiche.
\end{itemize}

Il confronto non mira a stabilire quale sistema sia superiore in assoluto, ma a evidenziare per quali classi di problemi ciascun paradigma risulti più adatto. Al fine di eliminare le alterazioni dovute alle operazioni di I/O a freddo su disco, ogni sessione prevede una fase iniziale di \textit{warm-up} finalizzata al completo riempimento delle rispettive Page Cache corporative.

\section{Scenario 1: La Potenza del Grafo (Neo4j vs RDBMS)}
\label{sec:scenario1}
Questo scenario analizza la capacità dei due motori di risolvere il cosiddetto \textit{Join Pain} quando la query naviga la topologia logica delle relazioni a profondità crescente.

\subsection{Query Multi-Hop (Analisi di Profondità)}
Il test prevede il calcolo ricorsivo della rete relazionale di un utente a profondità crescente (da 1 a 4 hop).
I risultati empirici rivelano un punto di inflessione critico tra il secondo e il terzo livello di profondità. A 1 hop, PostgreSQL risulta più veloce (impiega circa $0.33$ ms contro $0.76$ ms di Neo4j) grazie all'efficienza del lookup indicizzato in cache, mentre a 2 hop Neo4j passa in leggero vantaggio.

Il vero \textit{Join Pain} si manifesta chiaramente a partire dal terzo hop: Neo4j risponde in media in $\sim 3.6$ ms, mentre PostgreSQL subisce l'effetto del prodotto cartesiano generato dall'auto-join impiegando $\sim 45.9$ ms (speedup di $12.6\times$ a favore del grafo). Il limite strutturale del modello relazionale si palesa in modo drammatico al quarto hop: mentre Neo4j mantiene tempi di latenza ottimali e quasi costanti ($\sim 3.4$ ms), l'esplosione combinatoria porta il tempo di esecuzione di PostgreSQL a superare i $\sim 604$ ms (speedup empirico di oltre $178\times$). Questo conferma la superiorità del modello \textit{Index-Free Adjacency} per le navigazioni topologiche profonde.

\begin{table}[hbt!]
  \centering
  \begin{tabular}{lrrrr}
    \hline
    \textbf{Profondità} & \textbf{Neo4j (ms)} & \textbf{PostgreSQL (ms)} & \textbf{Speedup Neo4j} \\
    \hline
    1 Hop & 0.76 & 0.33 & 0.43$\times$ \\
    2 Hop & 1.14 & 1.48 & 1.30$\times$ \\
    3 Hop & 3.63 & 45.90 & 12.66$\times$ \\
    4 Hop & 3.38 & 604.88 & 178.80$\times$ \\
    \hline
  \end{tabular}
  \caption{Tempi medi di latenza per la query Multi-Hop.}
  \label{tab:multihop_results}
\end{table}

\begin{figure}[hbt!]
  \centering
  \includesvg[width=0.8\textwidth]{images/chapter4/scenario1/multihop_plot.svg}
  \caption{Esplosione combinatoria del tempo di esecuzione relazionale (scala logaritmica).}
  \label{fig:multihop_plot}
\end{figure}

Il limite del modello relazionale sulle query profonde non è puramente temporale, ma anche e soprattutto \textbf{spaziale}. Per dimostrarlo empiricamente, è stato condotto un monitoraggio in tempo reale dei container Docker isolando l'esecuzione della query a 4 hop. PostgreSQL, nel materializzare i self-join nidificati della CTE, ha richiesto dinamicamente al sistema operativo \textbf{+8.70 MB} di RAM aggiuntiva. Sebbene in termini assoluti possa sembrare un valore modesto, risulta in realtà spropositato se si considera che è l'allocazione \textit{dinamica per una singola query} effettuata su un dataset "giocattolo" di appena 1700 nodi. Su dataset a scala reale (milioni di nodi), questo comportamento moltiplicativo porta rapidamente la CTE a saturare la memoria disponibile, causando \textit{disk spilling} e azzerando le prestazioni. Al contrario, Neo4j pre-alloca la propria memoria tramite la JVM (motivo del suo baseline elevato a $\sim 5$ GB); le fluttuazioni registrate (circa +2 MB) non rappresentano allocazioni per dati intermedi della query, ma le normali e fisiologiche operazioni di \textit{Garbage Collection} e riorganizzazione dello heap Java in background, poiché la navigazione dei puntatori del grafo avviene \textit{in-place} senza duplicazione cartesiana dei dati.

\begin{table}[hbt!]
  \centering
  \begin{tabular}{lrr}
    \hline
    \textbf{Metrica RAM (SF 0.1)} & \textbf{PostgreSQL} & \textbf{Neo4j} \\
    \hline
    Baseline Iniziale & 34.83 MB & 5059.58 MB \\
    Picco Massimo & 43.53 MB & 5061.63 MB \\
    \textbf{Variazione Netta (Delta)} & \textbf{+8.70 MB} & \textbf{+2.05 MB} \\
    \hline
  \end{tabular}
  \caption{Costo spaziale in RAM per la risoluzione di un cammino a 4 hop.}
  \label{tab:ram_usage}
\end{table}

\begin{figure}[hbt!]
  \centering
  \includesvg[width=0.8\textwidth]{images/chapter4/scenario1/ram_usage_plot.svg}
  \caption{Profilazione dinamica della variazione di RAM (costo spaziale della CTE vs Grafo).}
  \label{fig:ram_plot}
\end{figure}

\subsection{Pattern Matching e Ricerca di Cicli}
Viene testata la capacità di identificare strutture ad anello chiuso all'interno del grafo delle transazioni o delle amicizie (strutture del tipo $A \rightarrow B \rightarrow C \rightarrow A$), pattern tipico nei modelli di \textit{Fraud Detection}.

In questo test, che rappresenta una query statistica disconnessa dalla topologia mirata all'intero database, l'RDBMS si dimostra superiore. Il conteggio globale dei triangoli favorisce PostgreSQL con uno speedup di oltre $17\times$ (impiegando $\sim 46$ ms contro gli $\sim 812$ ms di Neo4j). La ragione è strutturale: il motore SQL sfrutta algoritmi di \textit{hash join} altamente ottimizzati e le statistiche di colonna sull'intera tabella per operare a blocchi, mentre il database graph-native è costretto a materializzare il pattern matching esplorando l'intero grafo, non potendo beneficiare di filtri di tipo \textit{push-down} efficaci.

\begin{figure}[hbt!]
  \centering
  \includesvg[width=0.6\textwidth]{images/chapter4/scenario1/triangle_plot.svg}
  \caption{Confronto prestazioni nell'aggregazione statistica globale dei triangoli.}
  \label{fig:triangle_plot}
\end{figure}

\subsection{Pathfinding (Cammini Minimi)}
L'esperimento misura l'efficienza nel calcolare il percorso più breve tra due entità casuali nella rete sociale. Affidando il calcolo pathfinding interamente ai rispettivi motori di database, PostgreSQL è stato istruito per utilizzare una complessa \textit{CTE ricorsiva} dotata di un array ausiliario per la prevenzione dei cicli infiniti, mentre Neo4j ha impiegato il suo algoritmo nativo \texttt{shortestPath}.

I risultati mostrano una supremazia schiacciante da parte del database a grafo: Neo4j completa l'operazione in $\sim 0.9$ ms contro i $\sim 35$ ms di PostgreSQL. La topologia \textit{Index-Free Adjacency} e le euristiche di navigazione interne consentono a Neo4j di calcolare i cammini minimi con uno speedup sistematico di circa $40\times$ rispetto al massiccio aggravio computazionale introdotto dai continui auto-join della CTE relazionale, che satura rapidamente le proprie strutture dati ad ogni step di profondità.

\subsection{Sintesi Architetturale}
I risultati di questo primo scenario dimostrano come la scelta del database dipenda strettamente dalla classe di interrogazione. PostgreSQL eccelle nelle operazioni statistiche globali grazie all'efficienza degli hash join e all'elaborazione vettoriale. Tuttavia, nelle navigazioni topologiche profonde, Neo4j si impone categoricamente neutralizzando il fenomeno del \textit{Join Pain} sia in termini di latenza (speedup $>178\times$) sia per quanto concerne l'allocazione incontrollata di memoria.

A margine delle metriche prestazionali, l'architettura a grafo nativo offre due vantaggi strategici. Il primo riguarda la \textbf{flessibilità dello schema}: l'inserimento di nuove semantiche di relazione (es. "SEGUE" o "COLLABORA\_CON") in Neo4j richiede una semplice etichettatura degli archi, laddove PostgreSQL imporrebbe gravose migrazioni strutturali e la creazione di nuove tabelle associative. Il secondo vantaggio risiede nelle \textbf{potenzialità analitiche}: strutture topologiche complesse abilitano algoritmi di \textit{Graph Data Science} (come PageRank o Community Detection) direttamente all'interno del database, operazioni che nel mondo relazionale richiederebbero pesanti pipeline ETL verso framework di calcolo distribuiti esterni.

\section{Scenario 2: Transazioni e Concorrenza (Stress Test)}
\label{sec:scenario2}
Il secondo scenario mette sotto sforzo la robustezza transazionale di Neo4j in ambienti multi-utente ad alta concorrenza, focalizzandosi sul livello di isolamento \textit{Read Committed}. Lo stress test ha impiegato 8 thread concorrenti su un nodo \texttt{Person} condiviso, eseguendo 20 ripetizioni per ciascun sotto-test.

\subsection{Gestione del Livello Read Committed}
Viene verificata la stabilità di Neo4j durante l'esecuzione simultanea di 4 thread di lettura analitica e 4 thread di scrittura concorrenti sul medesimo nodo (proprietà \texttt{notification\_count}). Per rendere il test metodologicamente rigoroso, ogni thread scrittore utilizza una \textbf{transazione esplicita lunga}: prima scrive un valore non committato (\texttt{dirty\_value = 42}), aspetta 50--100 ms mantenendo la transazione aperta, poi aggiorna il valore finale (\texttt{committed\_value = 99}) e committa. I thread di lettura girano concorrentemente durante questa finestra: in un sistema senza isolamento vedrebbero il valore 42; in un sistema con \textit{Read Committed} possono osservare esclusivamente 0 (pre-scrittura) o 99 (post-commit).

Su un totale di 400 letture effettuate, \textbf{zero anomalie di tipo \textit{Dirty Read} sono state rilevate}: nessun reader ha mai osservato il valore 42 non committato, confermando il corretto funzionamento dei lock esclusivi di scrittura (X-lock) che Neo4j mantiene aperti per l'intera durata della transazione.

\begin{table}[hbt!]
  \centering
  \begin{tabular}{lrr}
    \hline
    \textbf{Metrica} & \textbf{Lettura} & \textbf{Scrittura} \\
    \hline
    Dirty Read rilevati & \multicolumn{2}{c}{0 $\checkmark$} \\
    Letture totali effettuate & \multicolumn{2}{c}{640} \\
    Latenza media (ms) & 1.798 & 265.082$^{*}$ \\
    Latenza mediana (ms) & 1.431 & 201.169 \\
    Latenza P90 (ms) & 2.796 & 546.069 \\
    \hline
  \end{tabular}
  \caption{Latenze Read Committed con 8 thread concorrenti (4 reader + 4 writer). $^{*}$La latenza del writer include la pausa esplicita di 50--100 ms che mantiene aperta la finestra di Dirty Read.}
  \label{tab:read_committed}
\end{table}



\subsection{Simulazione del Lost Update}
Viene forzata una \textit{race condition} applicativa in cui 10 thread concorrenti tentano di incrementare la medesima proprietà numerica (\texttt{notification\_count}) di un nodo \texttt{Person}, ciascuno leggendo lo stesso valore stale e riscrivendo il proprio risultato. Valore atteso a fine corsa: 10.

I risultati rivelano una distinzione netta tra le due strategie:

\begin{itemize}
  \item \textbf{Strategia non atomica} (\texttt{READ value; SET value = old + 1}): su 20 trial, \textbf{tutti e 20 producono lost update} (20/20 trial con perdita), con una media di \textbf{8.95 aggiornamenti persi per trial}. In casi estremi il valore finale osservato è 1 invece di 10, evidenziando che il livello \textit{Read Committed} non protegge da questo pattern poiché non acquisisce lock esclusivi automatici sulla sola lettura.
  \item \textbf{Strategia atomica} (\texttt{SET n.notification\_count = n.notification\_count + 1}): su 20 trial, \textbf{zero lost update} (0/20 trial con perdita). Il compilatore Cypher rileva la dipendenza self-referenziale e acquisisce preventivamente un lock esclusivo di scrittura (X-lock) sul nodo prima della lettura, rendendo l'intera operazione atomica.
\end{itemize}

\begin{table}[hbt!]
  \centering
  \begin{tabular}{lrrl}
    \hline
    \textbf{Strategia} & \textbf{Media Lost Update/Trial} & \textbf{Trial con perdita} & \textbf{Correttezza} \\
    \hline
    Non Atomica (alias separato) & 9.00 & 20/20 & $\times$ \\
    Atomica (\texttt{SET} diretto) & \textbf{0.00} & \textbf{0/20} & \checkmark \\
    \hline
  \end{tabular}
  \caption{Confronto tra strategia non atomica e atomica nella prevenzione del Lost Update.}
  \label{tab:lost_update}
\end{table}



\subsection{Deadlock Detection e Risoluzione}
Attraverso una coppia di thread con scritture incrociate simmetriche su nodi condivisi (Thread~A: acquisisce lock su Nodo~A poi tenta Nodo~B; Thread~B: acquisisce lock su Nodo~B poi tenta Nodo~A), viene forzata l'insorgenza deterministica di un'attesa circolare. Il handshake a due eventi garantisce che entrambe le transazioni tengano aperto il proprio lock simultaneamente prima di tentare il secondo.

Il \textit{Lock Manager} di Neo4j, tramite il monitoraggio continuo del \textit{Wait-for Graph} (grafo delle attese), intercetta il ciclo e interrompe la transazione causale lanciando una \texttt{TransientError} (codice \texttt{Neo.TransientError.Transaction.DeadlockDetected}) in un tempo medio di \textbf{35.53 ms} (P90: 71.49 ms), misurato dall'istante in cui il secondo lock viene richiesto fino alla ricezione dell'eccezione.

\begin{table}[hbt!]
  \centering
  \begin{tabular}{lr}
    \hline
    \textbf{Metrica} & \textbf{Valore} \\
    \hline
    Deadlock rilevati & 20 / 20 run \\
    Rollback automatici & 20 / 20 \checkmark \\
    Tempo di rilevazione medio (ms) & 35.53$^{*}$ \\
    Tempo di rilevazione P90 (ms) & 71.49$^{*}$ \\
    Meccanismo & Wait-for Graph \\
    Eccezione & \texttt{TransientError} \\
    \hline
  \end{tabular}
  \caption{Metriche di deadlock detection. $^{*}$Tempo misurato dall'istante in cui il secondo lock \`e richiesto fino all'arrivo della \texttt{TransientError}: misura l'overhead puro del Wait-for Graph, escludendo il delay artificiale di sincronizzazione.}
  \label{tab:deadlock}
\end{table}



\subsection{Sintesi Architetturale}
I risultati dello Scenario 2 confermano la solidità del modello transazionale di Neo4j in ambienti ad alta concorrenza. Il livello \textit{Read Committed} garantisce l'assenza di \textit{Dirty Read} in qualsiasi regime di contesa. Tuttavia, come dimostrato empiricamente, tale livello di isolamento \textbf{non è sufficiente a prevenire il \textit{Lost Update}} nel pattern read-modify-write separato: la soluzione nativa di Cypher (operazione self-referenziale atomica) è la strategia raccomandata per operazioni di incremento concorrente. Infine, il meccanismo di \textit{deadlock detection} basato sul Wait-for Graph si rivela robusto e trasparente: la risoluzione automatica tramite rollback avviene in tempi dell'ordine di 35--75 ms, senza necessità di timeout esterni a livello applicativo.

\section{Scenario 3: Sistemi Distribuiti e Teorema CAP}
\label{sec:scenario3}
Lo scenario sposta l'orizzonte verso le architetture aziendali distribuite ad alta affidabilità e scalabilità orizzontale. La sperimentazione è condotta su un cluster Neo4j 5.20.0 Enterprise con topologia a cinque nodi: tre \textit{Database Primaries} (partecipanti al consenso Raft) e due \textit{Database Secondaries} (repliche asincrone in sola lettura). Il dataset è LDBC SF~0.1 (1.700 nodi \texttt{Person}, 18.135 archi \texttt{KNOWS}), caricato via \texttt{LOAD CSV} dopo la formazione del quorum.

\subsection{Configurazione e Verifica del Cluster}
Il cluster viene avviato tramite \texttt{docker compose} con volumi Docker distinti per ciascun nodo, condizione necessaria affinché Neo4j assegni un \texttt{serverId} univoco a ogni istanza. L'utilizzo di volumi condivisi o copiati da un'istanza standalone causa un conflitto di identità Raft che impedisce la formazione del quorum.

Il Leader viene identificato interrogando la vista di sistema \texttt{SHOW DATABASES WHERE name = 'neo4j' AND writer = true} — l'API \texttt{dbms.cluster.overview()}, disponibile nelle versioni precedenti, è stata rimossa in Neo4j~5. Il cluster raggiunge il quorum e risponde alle prime query entro \textbf{5 secondi} dall'avvio dei container.

\begin{table}[hbt!]
  \centering
  \begin{tabular}{lll}
    \hline
    \textbf{Container} & \textbf{Ruolo Raft} & \textbf{Stato DBMS} \\
    \hline
    neo4j-core1      & Primary – Leader (writer) & Enabled \\
    neo4j-core2      & Primary – Follower        & Free    \\
    neo4j-core3      & Primary – Follower        & Enabled \\
    neo4j-secondary1 & Secondary                 & Free    \\
    neo4j-secondary2 & Secondary                 & Enabled \\
    \hline
  \end{tabular}
  \caption{Topologia del cluster al momento del benchmark. Lo stato \textit{Free} indica nodi registrati nella topologia DBMS ma privi di allocazione del database \texttt{neo4j}: il sistema di \textit{Autonomous Clustering} ha distribuito automaticamente tre repliche (core1, core3, secondary2), lasciando gli altri nodi come riserva.}
  \label{tab:cluster_topology}
\end{table}

\subsection{Tolleranza ai Guasti (Fault Tolerance)}
Durante un flusso continuo di transazioni in scrittura (4 thread paralleli per 15 secondi), viene simulato un crash hardware brutale tramite lo spegnimento forzato del container Leader (segnale \texttt{SIGKILL}, tempo di stop = 0). Ogni transazione è registrata con timestamp assoluto e latenza, consentendo di identificare con precisione il primo errore successivo al crash.

Il sistema risponde in linea con il profilo \textbf{CP} del Teorema CAP: le scritture vengono interrotte \textbf{8 ms} dopo il crash, e il cluster non le riprende finché non viene eletto un nuovo leader con quorum assoluto.

\begin{table}[hbt!]
  \centering
  \begin{tabular}{lr}
    \hline
    \textbf{Metrica} & \textbf{Valore} \\
    \hline
    Scritture totali registrate          & 1.047 \\
    Scritture riuscite                   & 1.001 (95,6\%) \\
    Scritture fallite (\textit{lost writes}) & 46 (4,4\%) \\
    Primo errore post-crash              & \textbf{182,7 ms} \\
    Downtime misurato scritture          & \textbf{24,34 s} \\
    Rielezione Raft completata (max 60s) & $\checkmark$ Sì \\
    Partition Tolerance sulle letture    & 8.961 / 8.961 \\
    Profilo CAP validato                 & \textbf{CP} \\
    \hline
  \end{tabular}
  \caption{Metriche del test di fault tolerance. Il cluster ha completato l'elezione di un nuovo leader in circa 24 secondi. Durante la finestra di downtime, il cluster ha sospeso le scritture per preservare la consistenza, ma ha continuato a processare con successo il 100\% delle letture (Partition Tolerance), garantendo il profilo CP perfetto.}
  \label{tab:fault_tolerance}
\end{table}

\begin{table}[hbt!]
  \centering
  \begin{tabular}{lrr}
    \hline
    \textbf{Statistica latenza} & \textbf{Scritture OK (ms)} & \textbf{Scritture fallite (ms)} \\
    \hline
    Media   & 29,6  & 190,6  \\
    Mediana & 7,2   & 107,8  \\
    P90     & 83,4  & 134,6  \\
    Min     & 2,2   & 1,4    \\
    Max     & 1.918,2 & 30.125,4 \\
    \hline
  \end{tabular}
  \caption{Distribuzione della latenza prima (scritture riuscite) e dopo (scritture fallite) il crash del Leader. Il valore massimo di 30.125~ms per le scritture fallite corrisponde al timeout di attesa della sonda.}
  \label{tab:fault_latency}
\end{table}

Il comportamento osservato conferma la scelta progettuale di Neo4j: \textbf{in assenza di quorum, nessuna scrittura viene accettata}. Questa garanzia elimina il rischio di \textit{split-brain} — scenario in cui due nodi si credono entrambi leader e accettano scritture divergenti — a costo di una temporanea indisponibilità controllata. In un deployment di produzione con allocazione esplicita del database su tutti i nodi e parametri Raft ottimizzati, il failover si completa tipicamente entro 10–30 secondi.

\subsection{Causal Consistency via Bookmark}
Neo4j 5 garantisce la \textit{Causal Consistency} attraverso un meccanismo di \textit{Bookmark}: al termine di una transazione in scrittura, il driver riceve un token che identifica la posizione nel log Raft. Le successive sessioni in lettura, inizializzate con quel Bookmark, vengono sospese dal server finché la replica destinataria non ha applicato il log-shipping fino al punto indicato — garantendo che il lettore veda sempre le proprie scritture, anche se instradata su un secondario asincrono.

Il test esegue una scrittura della proprietà \texttt{causal\_marker} su un nodo \texttt{Person}, recupera il Bookmark, e verifica 30 letture consecutive su nodi READ con il Bookmark allegato alla sessione.

\begin{table}[hbt!]
  \centering
  \begin{tabular}{lr}
    \hline
    \textbf{Metrica} & \textbf{Valore} \\
    \hline
    Letture corrette (valore aggiornato) & \textbf{30 / 30} \\
    Stale read osservate                 & \textbf{0} \\
    Overhead medio Bookmark              & 59,7 ms \\
    Overhead mediano Bookmark            & 3,2 ms \\
    P90                                  & 177,8 ms \\
    \hline
  \end{tabular}
  \caption{Risultati del test di Causal Consistency. Zero stale read su 30 letture: il Bookmark garantisce sempre la visibilità della propria scrittura su repliche asincrone. L'overhead mediano è di appena 3~ms, mentre la media riflette occasionali attese fisiologiche per l'allineamento del log sulle repliche.}
  \label{tab:causal_consistency}
\end{table}

Il risultato è ottimale: nessuna stale read in 30 letture. L'overhead mediano del Bookmark risulta irrisorio, configurando il meccanismo come una soluzione pratica per garantire la consistenza causale senza dover puntare sempre al leader — operazione che rappresenterebbe un collo di bottiglia sulla scalabilità in lettura.

\subsection{Scalabilità in Lettura e Load Balancing}
Tramite le funzionalità di cluster routing avanzato su protocollo Bolt nativo, è stato testato un pesante carico puramente in lettura con 16 thread paralleli. Le query sono state indirizzate uniformemente ai nodi abilitati per i carichi di lettura.

\begin{table}[hbt!]
  \centering
  \begin{tabular}{lr}
    \hline
    \textbf{Metrica} & \textbf{Valore} \\
    \hline
    Query totali registrate & 8.594 \\
    Query riuscite & 8.594 \\
    Throughput (\textit{QPS}) & \textbf{429,7} \\
    Latenza media & 37,3 ms \\
    Latenza P90 & 77,1 ms \\
    \hline
  \end{tabular}
  \caption{Risultati del test di scalabilità in lettura. Il sistema gestisce 430 query al secondo ripartendo in modo trasparente il carico su tutti i server della topologia al netto del Leader.}
  \label{tab:read_scalability}
\end{table}

Il driver Neo4j distribuisce efficacemente il traffico, confermando una eccellente scalabilità orizzontale in lettura del cluster.

\subsection{Sintesi Architetturale}
I risultati dello Scenario~3 validano empiricamente il posizionamento di Neo4j~5 Enterprise come sistema \textbf{CP} nel Teorema CAP:

\begin{itemize}
  \item \textbf{Consistency} $\checkmark$: nessuna scrittura incoerente viene mai accettata; il cluster sospende le operazioni in assenza di quorum.
  \item \textbf{Partition Tolerance} $\checkmark$: il protocollo Raft garantisce la coerenza del log anche in presenza di partizioni di rete.
  \item \textbf{Availability} (condizionale): le letture rimangono disponibili tramite il routing \texttt{neo4j://}; le scritture diventano temporaneamente indisponibili durante il failover del leader.
\end{itemize}

La Causal Consistency via Bookmark si dimostra robusta e a basso overhead, configurando Neo4j come soluzione adatta a contesti in cui la coerenza dei dati è non negoziabile e si può tollerare una breve indisponibilità controllata durante gli eventi di failover.



\section{Scenario 4: I Punti Deboli (Quando NON usare Neo4j)}
\label{sec:scenario4}
In ottica ingegneristica, questo paragrafo analizza i limiti strutturali del paradigma graph-native, identificando i contesti in cui l'architettura tabellare relazionale risulta nettamente superiore.

\subsection{Full-Table Scan e Aggregazioni Globali}
Viene eseguita una query puramente statistica disconnessa dalla topologia del grafo: il calcolo della lunghezza media dei testi di tutti i messaggi presenti nel database, raggruppati per tipologia di browser utilizzato.
\begin{itemize}
  \item \textbf{PostgreSQL} esegue un'operazione di \textit{Sequential Scan} parallelo sui blocchi contigui di memoria del disco, completando l'aggregazione di massa in frazioni di secondo.
  \item \textbf{Neo4j} evidenzia tempi di risposta significativamente più alti. Non disponendo di un motore colonnare o di strutture sequenziali piatte, il sistema deve scorrere i nodi sparsi nel supporto di memorizzazione per estrarre le proprietà, generando continui fenomeni di \textit{cache miss}.
\end{itemize}

\subsection{Inserimento Massivo di Dati Disconnessi (Bulk Insert)}
Il test misura il tempo richiesto per l'ingestione di un elevato volume di record grezzi privi di relazioni immediate (es. l'importazione di un'anagrafica piatta da un file CSV con un milione di righe). PostgreSQL, sfruttando l'efficienza del comando nativo \texttt{COPY}, surclassa nettamente Neo4j. Quest'ultimo soffre l'overhead computazionale legato all'architettura a grafo, la quale tenta di allocare strutture di memoria per i puntatori dei record anche in assenza di archi logici attivi.

\subsection{Esplosione Combinatoria nei Cammini Non Filtrati}
L'ultimo benchmark esplora i rischi legati a query Cypher mal ottimizzate o prive di vincoli di profondità espliciti su nodi ad alto branching factor (super-nodi). L'esecuzione di una query con cammini a lunghezza indefinita del tipo \texttt{-[*1..6]-} provoca la crescita esponenziale della frontiera di esplorazione. Nei log di sistema viene registrato il rapido consumo della memoria RAM del server, dimostrando come l'assenza di filtri topologici possa portare alla saturazione delle risorse hardware (\textit{Out Of Memory}) causata dalla memorizzazione di milioni di percorsi intermedi.
